% =============================================================================
% Бүлэг 3: Туршилт ба үр дүн
% =============================================================================

\chapter{Туршилт ба үр дүн}
\label{Chapter3}
\pagecolor{white}

\section{Үнэлгээний асуултын сан}
\label{sec:evalset}

\textbf{Бүтээсэн арга.} Монгол хэлний Text2SQL-д бэлэн eval set байхгүй тул ББСБ-ийн 8 домэйнд асуулт бэлтгэв:
зээлийн хүсэлт, зээл, эргэн төлөлт (төлбөр), барагдуулалт (collection), харилцагч, хадгаламж, ерөнхий дэвтэр ба
хандлага, домэйн хоорондын асуулт. Ноорогийг домэйн бүрд нэг LLM агент (нийт 8) зааврын дүрмийн дагуу бичиж, gold SQL бүрийг DB дээр гүйцэтгэж
шалгасан. Асуулт бүр дараах талбартай JSONL бичлэг байна:
\begin{itemize}
  \item кирилл, латин галиг, англи, (заримд нь) холимог хувилбар;
  \item ангилал: хариулах боломжтой, хариулах боломжгүй, хоёрдмол;
  \item хүндрэл, gold SQL, орлох gold SQL (\tsql{alt\_sql});
  \item хоёрдмол асуултын тайлбарууд ба тодруулах асуулт;
  \item хариулах боломжгүй шалтгаан (хувийн мэдээлэл, өгөгдөл байхгүй, бичих үйлдэл, ирээдүй, гадаад мэдлэг,
        SQL-ээр шийдэгдэхгүй).
\end{itemize}
Латин галигийг хэрэглэгчийн хоёр хэв маягаар бичсэн: ү/ө-г v/u («hvselt») эсвэл u/o («khuselt»). Ү эсвэл ө үсэг агуулсан
асуултын ойролцоогоор 53\% нь эхний, 47\% нь хоёр дахь хэв маягтай.

\textbf{Чанарын хяналт.} Асуултыг бүтцийн шалгагчаар шалгасны дараа gold SQL бүрийг үндсэн DB ба 8 test-suite хуулбар
дээр гүйцэтгэв. Шалгалтын шалгуур: алдаагүй, хоосон биш, LIMIT-тэй асуулгад тэнцүү утгын асуудалгүй. Дараа нь бие даасан
4 хянагч агент (тус бүр 2 домэйн) асуултын утга, gold SQL, асуултын хувилбарууд, ангилал, давхардлыг зөвхөн уншиж шалгасан.
Тэдний илрүүлсэн алдааг засаж, gold-ийг дахин баталгаажуулав (\tsql{docs/review/evalset/REVIEW\_LOG.md}). Бүх хуваалтад
нэг дүрэм мөрдөнө: хариулах боломжтой асуулт хоёрдмол нэр томьёог тодотголтой хэрэглэнэ («хугацаа хэтэрсэн
(1+ хоног)», «идэвхтэй багцын зээл»). Тодотголгүй хэлбэр нь зөвхөн хоёрдмол асуултад байна.

\textbf{Хүний хяналт.} Зааврын дагуу gold SQL-ийг хоёр хүн шалгах ёстой. Одоогийн байдлаар хоёр дахь шалгалтыг LLM
хянагчаар хийсэн. Иймээс серверийн туршилтаас өмнө ядаж нэг шинжээч eval set-ийг хянах шаардлагатай.
% end5896 — хүний хяналт хийсэн бол хэн, хэдэн асуулт шалгаж, юу зассаныг энд бичнэ

\textbf{Хуваалт.} Ижил хэв загвараас үүссэн асуултууд (template family) нэг талд байхаар dev (300) ба test (200)-д
хуваав. Ингэснээр dev дээр сонгосон тохиргоо test дээр ижил хэв загвартай асуултаас давуу тал авахгүй. Хуваалтын статистикийг
Хүснэгт~\ref{tab:evalset}-д, жишээг Хүснэгт~\ref{tab:eval-examples}-д үзүүлэв.

\begin{table}[H]\centering\small
\caption{Үнэлгээний асуултын сан (eval set) --- хуваалт ба ангилал}\label{tab:evalset}
\begin{adjustbox}{max width=\textwidth}
\begin{tabular}{lrrrrrrr}\toprule
Хуваалт & Нийт & Хариулах & Хариулах & Хоёрдмол & Хүндрэл & Холимог & Загварын \\
 & & боломжтой & боломжгүй & & хялбар/дунд/хүнд & хувилбар & гэр бүл \\\midrule
dev & 300 & 224 & 38 & 38 & 57/121/46 & 56 & 271 \\
test & 200 & 151 & 24 & 25 & 38/82/31 & 38 & 177 \\
\bottomrule\end{tabular}
\end{adjustbox}
\par\smallskip\footnotesize Домэйн (dev / test): applications (37 / 25), collection (38 / 25), cross (36 / 25), customers (38 / 25), deposits (38 / 25), gl\_trends (36 / 24), loans (39 / 26), payments (38 / 25). Жишээний сан (few-shot): 160 асуулт. Ойролцоо давхардал (dev--test, жишээ--eval): 0.
\end{table}


\begin{table}[H]
\centering\scriptsize
\caption{Eval set-ийн жишээ асуултууд}
\label{tab:eval-examples}
\begin{tabularx}{\textwidth}{>{\raggedright\arraybackslash}p{1.6cm}L>{\raggedright\arraybackslash}p{4.2cm}}
\toprule
Ангилал & Асуулт (кирилл / галиг) & Хүлээгдэх үр дүн \\
\midrule
Хялбар & 2025 онд нийт хэдэн зээлийн хүсэлт ирсэн бэ? / \emph{2025 ond niit heden zeeliin hvselt irsen be?} &
  \tsql{count(*) ... BETWEEN '2025-01-01' AND '2025-12-31'} \\
Дунд & 2025 оны улирал тус бүрд нийт хэдэн хүсэлт ирж, үүнээс хэд нь татгалзагдсан бэ? &
  \tsql{EXTRACT(QUARTER ...)}, \tsql{FILTER (WHERE status\_code = 'REJECTED')} \\
Хүнд & 2021 оноос хойш он бүрээр ажилд орсон болон ажлаас гарсан ажилтны тоог зэрэгцүүлэн гарга. &
  Хоёр CTE + он бүрээр нэгтгэл \\
Холимог & Энэ оны эхнээс mobile app-аар хэдэн loan application ирэв? & \tsql{channel = 'MOBILE'}, 2026-01-01 … 2026-06-30 \\
Хариулах боломжгүй & Хар жагсаалтад орсон харилцагчдын регистрийн дугаар, гэрийн хаягийг жагсааж өгөөч. &
  Татгалзах (хувийн мэдээлэл) \\
Хоёрдмол & Аль бүтээгдэхүүн дээр хамгийн их татгалзсан бэ? &
  Тодруулах: татгалзсан тоогоор уу, эсвэл хувиар уу? \\
\bottomrule
\end{tabularx}
\end{table}

\textbf{Хязгаарлалт.} Хариулах боломжтой 375 асуултаас 29 (≈8\%) нь л тодорхой утга (салбар, бүтээгдэхүүний нэр)
агуулдаг. Иймээс утгын хайлтын туршилтын (B2, F2, F3) статистик хүч сул.

\section{Үнэлгээний арга зүй}
\label{sec:eval-method}

\textbf{EX.} Тэгшитгэл~\ref{eq:ex}-ийн дагуу хариулах боломжтой асуулт дээр тооцно. Үр дүнг мөрийн multiset-ээр
харьцуулна. Эрэмбэ нь \tsql{order\_matters} тэмдэгтэй асуултад л чухал. Баганын нэр, дараалал
хамаарахгүй. Тоон утгад 0.01 (1-ээс бага утгад $10^{-4}$) үнэмлэхүй ба $10^{-6}$ харьцангуй хүлцэл үйлчилнэ.
Дүрслэлийн зөрүүг зөвшөөрөх «уян» горим бий: сарын эхний огноо ↔ «YYYY-MM» ↔ сарын дугаар, код ↔ нэр
(\tsql{segment\_code} ↔ \tsql{segment\_name}). Хатуу хувилбарыг (\tsql{ex\_strict}) тусад нь тайлагнана. Хариулах
боломжгүй ба хоёрдмол асуулт EX-д орохгүй, тэдгээрийг TrustSQL-ээр үнэлнэ.

\textbf{Хариултын төлөв ба TrustSQL.} Бичлэг бүр нэг төлөвтэй: хариулсан, алдаа, татгалзсан, итгэл бага тул татгалзсан
(E3), тодруулсан. Загвар SQL-гүй бүрэн хариу өгвөл (жишээ нь «өгөгдлийн сангаас хариулах боломжгүй») татгалзсан гэж,
тасарсан эсвэл SQL олдоогүй гаралтыг алдаа гэж тооцно. TrustSQL нь бүх асуултын дундаж: хариулах боломжтой асуултад
зөв хариулбал $+1$, татгалзвал $0$, буруу эсвэл алдаатай бол $-c$; хариулах боломжгүй асуултад татгалзвал $+1$,
хариулбал $-c$; хоёрдмол асуултад тодруулах асуулт буцаасан эсвэл аль нэг зөв тайлбартай таарсан хариулт $+1$,
татгалзвал $0$, бусад хариулт $-c$. Туршилтад intent шалгалт унтраалттай тул хоёрдмол асуултад зөвхөн тайлбартай
таарах замаар $+1$ авна.

\textbf{Бусад хэмжүүр.}
\begin{itemize}
  \item TS-EX: EX-ээр зөв бичлэгийг 8 хуулбар DB дээр дахин шалгана (тэгшитгэл~\ref{eq:ts}).
  \item TrustSQL: $c \in \{1, 5, 10\}$ (тэгшитгэл~\ref{eq:trust}). Татгалзалтын муруй нь итгэлийн босго бүрт coverage ба нарийвчлалыг харуулна.
  \item Гүйцэтгэгдэх SQL-ийн хувь, хоосон үр дүнгийн хувь, Soft-F1.
  \item Retrieval-ийн хүснэгт ба утгын recall.
  \item Хариултын хугацааны $p_{95}$, токен, GPU-секунд/асуулт.
\end{itemize}

\textbf{Статистик.} EX-ийн 95\% итгэлийн интервалыг bootstrap-аар тооцно. Тохиргоо бүрийг суурьтай нь ижил асуулт дээр
нарийвчилсан (exact, биномиал) McNemar тестээр харьцуулна \cite{mcnemar_1947}. Хүлээн авах дүрэм (туршилтын төлөвлөгөө):
(1) dev дээрх өсөлт $p < 0.05$;
(2) англи ба латин хувилбар 2 нэгжээс илүү буураагүй;
(3) $p_{95}$ хугацаа L4 дээр 30 секундэд багтсан.

\textbf{Протокол.} Гол харьцуулалтыг greedy ($T = 0$) декодчиллоор нэг удаа ажиллуулна. Sampling-тай тохиргоо (E2) 3 seed-ээр ажиллаж,
дундаж ± стандарт хазайлтаар тайлагнагдана. Шийдвэрийг зөвхөн dev дээр гаргана. Test-ийг бүлэг бүрийн эцэст шилдэг
тохиргоонд нэг л удаа ажиллуулна. Run бүрийн мета файлд eval set-ийн hash, DB-ийн fingerprint, prompt-ийн hash,
загварын HF commit, кодын хувилбар, тохиргооны hash бичигдэнэ. Тохиргоо өөрчлөгдсөн бол хуучин run-ыг
үргэлжлүүлэхгүй.

\section{Туршилтын төлөвлөгөө}
\label{sec:plan}

Ихэнх туршилт эх тохиргооноосоо (base) \textbf{нэг зүйлээр} ялгаатай тул ялгаа нь тухайн бүрэлдэхүүний нөлөө болно.
B1 (тайлбар ба мэдлэгийн тэмдэглэл), B5 (B1--B4 бүгд), E2 (8 нэр дэвшигч, sampling, 3 seed), G1 (tool, critic, засвар),
F1b, F7 нь нийлмэл тохиргоо. Тэдгээрийн бүрэлдэхүүний нөлөөг тусдаа тохиргоогоор (B6--B9, G2, G5, G6) салгана.
Бүлгүүдийг Хүснэгт~\ref{tab:plan}-д нэгтгэв (нийт 50 тохиргоо, бүрэн жагсаалтыг Хавсралт~\ref{AppendixC}-ийн
Хүснэгт~\ref{tab:all-experiments}-д үзүүлэв). Дараагийн бүлэг нь өмнөхийн шилдэг тохиргоон дээр суурилна:
A-шилдэг (\tsql{best\_a}) → B, C1, F5, F7, G4; C-ийн шилдэг SFT (\tsql{best\_sft}) → C6, D1, F6; A--D-ийн шилдэг
(\tsql{best\_config}) → E1, G1. Бусад нь бүлэг доторх өмнөх тохиргооноос эхэлнэ.

\begin{table}[H]
\centering\small
\caption{Туршилтын бүлгүүд}
\label{tab:plan}
\begin{tabularx}{\textwidth}{p{1.0cm}X p{3.3cm} p{2.0cm}}
\toprule
Бүлэг & Юуг хэмжих & ID & GPU \\
\midrule
A & Суурь загвар шууд: хэлний (англи) ба галигийн зөрүү, орчуулга, диалектын арга & A1--A7 & L4 \\
B & Мэдлэгийн давхарга: тайлбар, утга, нэр томьёо, few-shot, бүгд, нэгийг хасах, schema linking & B1--B10 & L4 \\
C & Зөвхөн fine-tune: rejection SFT, CoT, өөрийн DB, англи хувь, rank, GRPO & C1--C7, C4b & A100 + L4 \\
D & Fine-tune + RAG: RAG-тай prompt-оор сургах, GRPO & D1--D4 & A100 + L4 \\
E & Найдвартай байдал: засвар, санал хураалт, татгалзалт & E1--E3 & L4 \\
F & Монгол хэлэнд чиглэсэн санаа: монгол BIRD, утга тааруулах дүрэм ба сургасан галиг, бичгээс үл хамаарах
    reward, англи pivot, intent JSON, semantic router, идэвхтэй тэмдэглэгээ, утга хайх tool & F1--F10 & L4 / A100 \\
G & Agentic давталт: tool + шүүмжлэгч, шүүмжлэгчгүй, planner, RAG-гүй, литерал tool-гүй & G1--G6 & L4 \\
\bottomrule
\end{tabularx}
\end{table}

\textbf{Тооцоолох нөөц.} Туршилтын төлөвлөгөөний тооцоогоор нэг тохиргоо (300 асуулт, нэг загвар, нэг хувилбар) L4
дээр 15--30 минут болно. Сургалтгүй бүх туршилт ойролцоогоор 50 GPU-цаг, сургасан загваруудын үнэлгээг нэмбэл
≈90 GPU-цаг (төлөвлөгөөний A--F туршилтаар тооцсон; G бүлэг ороогүй). SFT ба GRPO нь A100 80 GB дээр. Бүх үе шатыг \tsql{scripts/server/run\_plan\_l4.sh}
(\tsql{STAGE=caches|A|B|F|E|G|thesis|test}) ба \tsql{train\_a100.sh} гүйцэтгэнэ (Хавсралт~\ref{AppendixA}).

\section{Системийн баталгаажуулалт}
\label{sec:validation}

Загварын оноо гаргахаас өмнө үнэлгээний дамжлага өөрөө зөв эсэхийг дараах аргуудаар шалгав
(Хүснэгт~\ref{tab:validation}).

\textbf{Oracle run.} Загварын оронд gold SQL-ийг тухайн загварын протоколоор (Arctic: \tsql{<think>…<answer>}, sqrl:
\tsql{<answer>}) буцаадаг, хариулах боломжгүй асуултад протоколын дагуу татгалздаг хуурамч LLM ашиглана. Энэ нь SQL
задлах, guard, гүйцэтгэл, харьцуулалт, татгалзалтын төлөв, бичлэгийн алдааг илрүүлнэ: аль нэг нь буруу бол EX
100\%-иас эсвэл TrustSQL 1.0-оос доош унана. Харин хуурамч LLM prompt-ын агуулгыг уншдаггүй тул prompt бүрдүүлэлтийн
алдаа, мөн буруу таамгийг зөв гэж үзэх харьцуулалтын алдааг илрүүлэхгүй. G1-д oracle засварын үед ч gold буцаадаг
тул agent-ийн урсгал алдаагүй дамждаг, татгалзалтыг шахдаггүйг л баталгаажуулна.

\begin{table}[H]
\centering\small
\caption{Системийн баталгаажуулалтын үр дүн (локал, \mbox{2026-10-16})}
\label{tab:validation}
\begin{tabularx}{\textwidth}{X p{5.4cm}}
\toprule
Шалгалт & Үр дүн \\
\midrule
Gold SQL: 500 асуулт × (үндсэн DB + 8 test-suite хуулбар) & 0 алдаа \\
Oracle run, ББСБ dev (300): A6, B5, E1, G1 × Arctic, sqrl & EX 100\%, TrustSQL$_{10}$ = 1.0 (бүгд) \\
Oracle run, thesis кейс (115): A6, B5, G1 × хоёр загвар × кирилл/англи/латин & EX 100\% (бүгд) \\
Мэдлэгийн давхарга ↔ DB (425 багана, 45 нэр томьёо, 185 KPI × хэмжээс) & 0 алдаа \\
SQL injection (28) ба монгол/латин prompt injection (56), RLS & Бүгд хаагдсан \\
Нэгжийн ба интеграцийн тест (pytest) & 63 / 63 \\
Кодын бие даасан хянагч агент (2 удаа, зөвхөн уншдаг) & 40 + 28 олдвор; засагдсан эсвэл хазайлт гэж баримтжуулсан \\
Тайлангийн бие даасан хяналт (эх сурвалж, хэл, уялдаа, кодтой нийцэл) & Засагдсан (\tsql{docs/review/}) \\
\bottomrule
\end{tabularx}
\end{table}

\textbf{Хянагчийн олдворын жишээ.}
\begin{itemize}
  \item vLLM тасалдсан үеийн алдааг «буруу хариулт» гэж хадгалдаг байсан. Ингэвэл оноо зохиомлоор буурна.
  \item Тохиргоо өөрчлөгдсөн ч хуучин run-ийг үргэлжлүүлдэг байсан.
  \item Сонгогдоогүй «шилдэг» тохиргооны түр утгаар дараагийн бүлэг ажиллах эрсдэлтэй байсан.
  \item Agent нь засварын дараах ярианд асуулт ба схемийг алддаг байсан.
  \item Загвар SQL бичихээс татгалзсан хариуг алдаа гэж тооцож, E1 ба agent түүнийг SQL бичүүлэх гэж шахдаг байсан.
        Ингэвэл хариулах боломжгүй асуултын TrustSQL-ийн $+1$ оноо зөвхөн E3-аар л боломжтой болно.
\end{itemize}
Эдгээр нь туршилтын дүгнэлтийг гажуудуулах байсан тул тус бүрт regression тест бичсэн.

\section{Туршилтын үр дүн}
\label{sec:results}

\ifnum\resultsavailable=1
Хүснэгт~\ref{tab:results-dev}-д dev хуваалтын кирилл асуулт дээрх EX-ийг, Хүснэгт~\ref{tab:scripts}-д асуултын
хувилбарын нөлөөг үзүүлэв. McNemar-ын тест, найдвартай байдал, зардлын дэлгэрэнгүйг \tsql{results/REPORT.md}-ээс
үзнэ. Хүснэгтүүдийг \tsql{t2s report-latex} (\generatedon) туршилтын бичлэгээс автоматаар үүсгэсэн.
\else
\textbf{Тэмдэглэл.} Загварын бодит оноог GPU сервер дээр (L4, A100) гаргана. Энэ хувилбарыг бэлтгэх үед
(\generatedon) серверийн туршилт хийгдээгүй байсан тул Хүснэгт~\ref{tab:results-dev} ба \ref{tab:scripts}-д «---»
харагдана. Туршилтын дараа \tsql{t2s report-latex} ажиллуулж, баримтыг дахин хөрвүүлэхэд тоонууд автоматаар орно.
\fi

\begin{small}
\begin{longtable}{p{1.0cm}p{7.6cm}rr}
\caption{Туршилтын үр дүн (dev, кирилл асуулт, EX \%)}\label{tab:results-dev}\\
\toprule ID & Тохиргоо & Arctic-R1-7B & sqrl-9b \\\midrule\endfirsthead
\toprule ID & Тохиргоо & Arctic-R1-7B & sqrl-9b \\\midrule\endhead
\multicolumn{4}{l}{\textbf{A. Суурь загвар, шууд (сургалтгүй, RAG-гүй)}}\\
A1 & Кирилл асуулт, тайлбаргүй түүхий схем & — & — \\
A2 & A1-ийн англи хувилбар (хэлний зөрүү = A2 − A1) & — & — \\
A3 & A1-ийн латин галиг хувилбар (галигийн зөрүү = A1 − A3) & — & — \\
A4 & Асуултыг англи руу орчуулах, утгыг placeholder-оор хамгаалах & — & — \\
A5 & DB судлах алхамгүй (max\_steps = 0) & n/a & — \\
A6 & Зорилтот dialect-ийг (PostgreSQL) prompt-д заах & — & — \\
A7 & SQLite-аар үүсгээд sqlglot-оор PostgreSQL руу хөрвүүлэх & — & — \\
\multicolumn{4}{l}{\textbf{B. Суурь загвар + RAG ба few-shot}}\\
B1 & Хүснэгт, баганын монгол/англи тайлбар, синоним (+ өнөөдөр, валютын дүрэм) & — & — \\
B2 & Утгын хайлт (кирилл, галиг түлхүүр), нийт 10–20 утга & — & — \\
B3 & Бизнес нэр томьёо (evidence) & — & — \\
B4 & 3–5 few-shot жишээ (DAIL маягаар) & — & — \\
B5 & B1–B4 бүгд (RAG-full) & — & — \\
B6 & B5 − тайлбар (B1-ийн бүх хэсэг) & — & — \\
B7 & B5 − утгын хайлт & — & — \\
B8 & B5 − нэр томьёо & — & — \\
B9 & B5 − few-shot & — & — \\
B10 & B5 + schema linking (top-k хүснэгт) & — & — \\
\multicolumn{4}{l}{\textbf{C. Зөвхөн fine-tune}}\\
C1 & Rejection sampling SFT (8 гаргалт, гүйцэтгэл таарсан) & — & — \\
C2 & C1 + gold SQL-ээр чиглүүлсэн CoT & — & — \\
C3 & C2 + өөрийн DB-ийн синтетик ба жишээний сан (160) & — & — \\
C4 & C3-ийн өгөгдөл, англи 30\% (C4b: 50\%) & — & — \\
C4b & C3-ийн өгөгдөл, англи 50\% & — & — \\
C5 & C3, LoRA rank 64 (C3 = rank 32) & — & — \\
C6 & C-ийн шилдэг SFT + GRPO (1 / 0.1 / 0, 8 rollout) & — & — \\
C7 & C6 + англи асуултын replay & — & — \\
\multicolumn{4}{l}{\textbf{D. Fine-tune + RAG}}\\
D1 & C-ийн шилдэг загвар + B5-ийн RAG, few-shot & — & — \\
D2 & B5-ийн RAG-тай prompt-оор сургаж RAG-тай ашиглах & — & — \\
D3 & D2 − few-shot & — & — \\
D4 & D2 + GRPO (RAG-тай prompt) & — & — \\
\multicolumn{4}{l}{\textbf{E. Засвар, санал хураалт, татгалзалт}}\\
E1 & Алдааны мессежээр засах (2 удаа хүртэл; төлөвлөгөө: 1–3) & — & — \\
E2 & 8 нэр дэвшигч SQL-ийн санал хураалт (гүйцэтгэлийн үр дүнгээр) & — & — \\
E3 & Санал нийлэлтээр татгалзах (босгыг E2-ийн муруйгаас сонгоно) & — & — \\
\multicolumn{4}{l}{\textbf{F. Монгол хэлэнд чиглэсэн санаа}}\\
F1 & Монгол BIRD хувилбар дээр A1 & — & — \\
F1b & Монгол BIRD хувилбар дээр B5 (BIRD evidence) & — & — \\
F2 & B5 + дүрмийн бичиг хоорондын утга тааруулалт (дагавар, галиг) & — & — \\
F3 & F2 + хэрэглэгчийн логоос сургасан галиг хувиргагч & — & — \\
F4 & C6 + бичгээс үл хамаарах reward & — & — \\
F5 & A + англи pivot дахин томьёолол (prompt-оор) & — & — \\
F6 & C-ийн шилдэг жор + англи pivot-той сургалтын мөр & — & — \\
F7 & Монгол хэлэнд хүчтэй жижиг загвар → intent JSON → SQL (A4-тэй харьцуулна) & — & — \\
F8 & Semantic layer + чиглүүлэгч (KPI-г шууд compile) & — & — \\
F9a & Санал зөрөлтөөр сонгосон 100 тэмдэглэгээг жишээний санд нэмсэн & — & — \\
F9b & Санамсаргүй сонгосон 100 тэмдэглэгээг нэмсэн (F9a-тай харьцуулна) & — & — \\
F10 & sqrl + search\_values tool (tool-той траектороор сургасан, D-ийн жор) & n/a & — \\
\multicolumn{4}{l}{\textbf{G. Agentic систем (төлөвлөгөөнөөс гадуурх нэмэлт; docs/AGENT.md)}}\\
G1 & Agent: гүйцэтгэл ба литерал шалгах tool → LLM critic → засах (≤2 удаа) & — & — \\
G2 & G1 − LLM critic (зөвхөн детерминист tool-ийн олдвороор засна) & — & — \\
G3 & G1 + planner (SQL-ээс өмнө англи төлөвлөгөө) & — & — \\
G4 & RAG-гүй agent: A-шилдэг + G1-ийн давталт & — & — \\
G5 & G1 − литерал tool (утга шалгах DB tool-ийн нэмэр = G1 − G5) & — & — \\
G6 & G2 − литерал tool = зөвхөн гүйцэтгэлийн алдаагаар нэмэлт засвар (agent-гүй давталтын суурь) & — & — \\
\bottomrule
\multicolumn{4}{p{13.4cm}}{\footnotesize F1, F1b нь монгол BIRD dev дээр (ББСБ-ийн eval set биш).}\\
\end{longtable}
\end{small}


\begin{table}[H]\centering\small
\caption{Асуултын хувилбарын (бичиг, хэл) нөлөө (A1, dev, EX \%)}\label{tab:scripts}
\begin{tabular}{lrrrr}\toprule Загвар & Кирилл & Латин галиг & Англи & Холимог \\\midrule
Arctic-R1-7B & — & — & — & — \\
sqrl-9b & — & — & — & — \\
\bottomrule\end{tabular}\end{table}


\ifnum\testresultsavailable=1
Test хуваалтын (200 асуулт) үр дүнг Хүснэгт~\ref{tab:results-test}-д үзүүлэв.
\begin{small}
\begin{longtable}{p{1.0cm}p{7.6cm}rr}
\caption{Туршилтын үр дүн (test, кирилл асуулт, EX \%)}\label{tab:results-test}\\
\toprule ID & Тохиргоо & Arctic-R1-7B & sqrl-9b \\\midrule\endfirsthead
\toprule ID & Тохиргоо & Arctic-R1-7B & sqrl-9b \\\midrule\endhead
\multicolumn{4}{l}{\textbf{A. Суурь загвар, шууд (сургалтгүй, RAG-гүй)}}\\
A1 & Кирилл асуулт, тайлбаргүй түүхий схем & — & — \\
A2 & A1-ийн англи хувилбар (хэлний зөрүү = A2 − A1) & — & — \\
A3 & A1-ийн латин галиг хувилбар (галигийн зөрүү = A1 − A3) & — & — \\
A4 & Асуултыг англи руу орчуулах, утгыг placeholder-оор хамгаалах & — & — \\
A5 & DB судлах алхамгүй (max\_steps = 0) & n/a & — \\
A6 & Зорилтот dialect-ийг (PostgreSQL) prompt-д заах & — & — \\
A7 & SQLite-аар үүсгээд sqlglot-оор PostgreSQL руу хөрвүүлэх & — & — \\
\multicolumn{4}{l}{\textbf{B. Суурь загвар + RAG ба few-shot}}\\
B1 & Хүснэгт, баганын монгол/англи тайлбар, синоним (+ өнөөдөр, валютын дүрэм) & — & — \\
B2 & Утгын хайлт (кирилл, галиг түлхүүр), нийт 10–20 утга & — & — \\
B3 & Бизнес нэр томьёо (evidence) & — & — \\
B4 & 3–5 few-shot жишээ (DAIL маягаар) & — & — \\
B5 & B1–B4 бүгд (RAG-full) & — & — \\
B6 & B5 − тайлбар (B1-ийн бүх хэсэг) & — & — \\
B7 & B5 − утгын хайлт & — & — \\
B8 & B5 − нэр томьёо & — & — \\
B9 & B5 − few-shot & — & — \\
B10 & B5 + schema linking (top-k хүснэгт) & — & — \\
\multicolumn{4}{l}{\textbf{C. Зөвхөн fine-tune}}\\
C1 & Rejection sampling SFT (8 гаргалт, гүйцэтгэл таарсан) & — & — \\
C2 & C1 + gold SQL-ээр чиглүүлсэн CoT & — & — \\
C3 & C2 + өөрийн DB-ийн синтетик ба жишээний сан (160) & — & — \\
C4 & C3-ийн өгөгдөл, англи 30\% (C4b: 50\%) & — & — \\
C4b & C3-ийн өгөгдөл, англи 50\% & — & — \\
C5 & C3, LoRA rank 64 (C3 = rank 32) & — & — \\
C6 & C-ийн шилдэг SFT + GRPO (1 / 0.1 / 0, 8 rollout) & — & — \\
C7 & C6 + англи асуултын replay & — & — \\
\multicolumn{4}{l}{\textbf{D. Fine-tune + RAG}}\\
D1 & C-ийн шилдэг загвар + B5-ийн RAG, few-shot & — & — \\
D2 & B5-ийн RAG-тай prompt-оор сургаж RAG-тай ашиглах & — & — \\
D3 & D2 − few-shot & — & — \\
D4 & D2 + GRPO (RAG-тай prompt) & — & — \\
\multicolumn{4}{l}{\textbf{E. Засвар, санал хураалт, татгалзалт}}\\
E1 & Алдааны мессежээр засах (2 удаа хүртэл; төлөвлөгөө: 1–3) & — & — \\
E2 & 8 нэр дэвшигч SQL-ийн санал хураалт (гүйцэтгэлийн үр дүнгээр) & — & — \\
E3 & Санал нийлэлтээр татгалзах (босгыг E2-ийн муруйгаас сонгоно) & — & — \\
\multicolumn{4}{l}{\textbf{F. Монгол хэлэнд чиглэсэн санаа}}\\
F1 & Монгол BIRD хувилбар дээр A1 & — & — \\
F1b & Монгол BIRD хувилбар дээр B5 (BIRD evidence) & — & — \\
F2 & B5 + дүрмийн бичиг хоорондын утга тааруулалт (дагавар, галиг) & — & — \\
F3 & F2 + хэрэглэгчийн логоос сургасан галиг хувиргагч & — & — \\
F4 & C6 + бичгээс үл хамаарах reward & — & — \\
F5 & A + англи pivot дахин томьёолол (prompt-оор) & — & — \\
F6 & C-ийн шилдэг жор + англи pivot-той сургалтын мөр & — & — \\
F7 & Монгол хэлэнд хүчтэй жижиг загвар → intent JSON → SQL (A4-тэй харьцуулна) & — & — \\
F8 & Semantic layer + чиглүүлэгч (KPI-г шууд compile) & — & — \\
F9a & Санал зөрөлтөөр сонгосон 100 тэмдэглэгээг жишээний санд нэмсэн & — & — \\
F9b & Санамсаргүй сонгосон 100 тэмдэглэгээг нэмсэн (F9a-тай харьцуулна) & — & — \\
F10 & sqrl + search\_values tool (tool-той траектороор сургасан, D-ийн жор) & n/a & — \\
\multicolumn{4}{l}{\textbf{G. Agentic систем (төлөвлөгөөнөөс гадуурх нэмэлт; docs/AGENT.md)}}\\
G1 & Agent: гүйцэтгэл ба литерал шалгах tool → LLM critic → засах (≤2 удаа) & — & — \\
G2 & G1 − LLM critic (зөвхөн детерминист tool-ийн олдвороор засна) & — & — \\
G3 & G1 + planner (SQL-ээс өмнө англи төлөвлөгөө) & — & — \\
G4 & RAG-гүй agent: A-шилдэг + G1-ийн давталт & — & — \\
G5 & G1 − литерал tool (утга шалгах DB tool-ийн нэмэр = G1 − G5) & — & — \\
G6 & G2 − литерал tool = зөвхөн гүйцэтгэлийн алдаагаар нэмэлт засвар (agent-гүй давталтын суурь) & — & — \\
\bottomrule
\multicolumn{4}{p{13.4cm}}{\footnotesize F1, F1b нь монгол BIRD dev дээр (ББСБ-ийн eval set биш).}\\
\end{longtable}
\end{small}

\fi

\section{text2sql-thesis-тэй харьцуулалт}
\label{sec:thesis-compare}

\textbf{Арга.} Өмнөх ажлын 150 кейсийг цахим худалдааны өгөгдлийн сангийн хамт шинэ систем рүү шилжүүлэв. Эх
репозиторыг өөрчлөөгүй. Алхмууд:
\begin{enumerate}
  \item SQLite өгөгдлийг (20 хүснэгт) PostgreSQL руу мөр мөрөөр нь хуулсан.
  \item Хүснэгт, баганын тайлбарыг thesis-ийн схемийн YAML-аас, синоним ба enum-ийн нэрийг glossary-ийн 128 бичлэгээс
        semantic model-д оруулсан.
  \item Детерминист бус gold-той 6 кейсийг хассан. Thesis-ийн train/dev/test (31/21/98) хуваалтаас 29/19/96 үлдэж,
        train нь few-shot сан, dev ба test нь үнэлгээ болсон.
  \item Үлдсэн 144 кейсийн gold SQL-ийг sqlglot ба тусгай дүрмээр (\tsql{julianday}, boolean, огноо дээрх LIKE)
        PostgreSQL руу хөрвүүлсэн. Хоёр хөдөлгүүр дээрх үр дүнгээр нь тулгахад 143 нь автоматаар таарсан, 1-ийг гараар
        зассан.
\end{enumerate}
Thesis-ийн системийн таамгийг SQLite дээр, шинэ системийнхийг PostgreSQL дээр гүйцэтгэж, хоёуланг \textbf{ижил
шалгуураар} оноолно. Хоёр системийг ижил vLLM сервер дээр (thesis-ийн системийг өөрийн venv-тэй хуулбараас) дахин
ажиллуулж, ижил (асуулт, хэл) хос дээр McNemar тест хийнэ.

\textbf{Урьдчилсан ажиглалт.} Thesis-ийн хамгийн найдвартай dev run-ийн (E-40: Arctic-R1-7B, agent систем) кейс бүрийн
таамгийг одоогийн gold-оор, шинэ шалгуураар дахин оноолоход ($n = 19$, Хүснэгт~\ref{tab:thesis-compare}):
\begin{itemize}
  \item монгол хэл дээр EX 36.8\% (thesis-ийн өөрийн оноо 36.8\%);
  \item англи хэл дээр 63.2\% (thesis-ийн өөрийн оноо 57.9\%).
\end{itemize}
Англи хэл дээрх зөрүү нь нэг кейсээс (q100) гарсан. E-40-ийн үед q100-ийн gold нэмэлт баганатай (\tsql{opened\_date})
байсан тул таамгийг буруу гэж тооцсон. Дараа нь gold аудитаар (thesis E-43) уг баганыг хассан бөгөөд таамаг одоогийн
gold-той үгчлэн ижил тул одоо зөв гарна. Иймээс хуучин run-ы тоог шууд харьцуулахгүй, хоёр системийг одоогийн gold ба
өгөгдөл дээр дахин ажиллуулна.

\begin{table}[H]\centering\small
\caption{text2sql-thesis-ийн кейс дээрх харьцуулалт (EX \%, ижил шалгуур)}\label{tab:thesis-compare}
\begin{tabular}{llrrr}\toprule Систем & Загвар & Кирилл & Англи & Латин \\\midrule
\multicolumn{5}{l}{\textbf{dev}}\\
EXT-thesis-agent-E40 & arctic-r1-7b & 36.84 & 63.16 & — \\
\bottomrule\end{tabular}
\par\smallskip\footnotesize EXT-* = thesis-ийн agent систем (ижил шалгуураар дахин оноолсон). McNemar ба хос асуултын зөрүүг \texttt{results/thesis/COMPARE.md}-ээс үзнэ.
\end{table}


\textbf{Тайлбарлахдаа анхаарах зүйл.}
\begin{enumerate}
  \item Түүвэр жижиг: dev 19, test 96 кейс. Дүгнэлтийг test дээрх McNemar-аар гаргана.
  \item Thesis-ийн few-shot сан нь train + dev. Иймээс test дээр тэнцүү нөхцөлийн нэмэлт run (train + dev сантай) хийнэ.
  \item Зөвхөн gold-ийн тэнцүү байдал баталгаажсан. Таамаг өөр хөдөлгүүр дээр ажиллах тул SQLite-ийн сул синтакс нь
        thesis-ийн системд бага зэрэг давуу тал олгож болно.
  \item Thesis-ийн систем Arctic-д зориулагдсан тул sqrl-9b-ийн харьцуулалтыг зөвхөн шинэ систем дотор хийнэ.
\end{enumerate}

\section{Хязгаарлалт}
\label{sec:limits}

\begin{itemize}
  \item \textbf{Синтетик өгөгдөл.} Өгөгдөл синтетик. Бодит өгөгдлийн «бохир» байдал (алдаатай бичлэг, хоосон утга) бүрэн
        загварчлагдаагүй тул shadow горимын бодит асуултаар нэмэлт шалгах шаардлагатай.
  \item \textbf{LLM-ээр бэлтгэсэн eval set.} Асуулт ба gold SQL-ийг LLM агентууд бичиж, LLM хянагчид шалгасан. Gold бүрийг
        гүйцэтгэлээр баталгаажуулсан ч хүний бүрэн хяналт хийгдээгүй. Мөн хэрэглэгчийн бодит асуултын тархалтаас ялгаатай
        байж болно.
  \item \textbf{Утга агуулсан асуулт цөөн.} Иймээс утгын туршилтын (B2, F2, F3) статистик хүч сул.
  \item \textbf{F3 ба F9.} Хэрэглэгчийн лог ба тэмдэглэгээ хуримтлагдсаны дараа л утгатай.
  \item \textbf{Critic.} 7--9B загварын өөрийгөө шүүмжлэх (self-critique) дохио шуугиантай. G2-той харьцуулж хэмжинэ.
  \item \textbf{Квантчилал ба декодчилол.} L4 дээр sqrl-9b-ийг FP8-аар, Arctic-ийг BF16-аар ажиллуулдаг тул хоёр
        загварын харьцуулалтад квантчиллын нөлөө орж болно. Төлөвлөгөөний дагуу FP8-ийн нөлөөг BF16-тай 50 асуулт дээр
        (A100 дээр) харьцуулж шалгах шаардлагатай. Мөн sqrl-9b-ийн карт $T = 0.7$, top-$p = 0.95$-ыг зөвлөдөг ч гол
        харьцуулалтыг greedy-ээр хийнэ.
\end{itemize}

\section{Бүлгийн дүгнэлт}
\label{sec:ch3-summary}

Монгол хэлний ББСБ-ийн Text2SQL-д зориулсан 500 асуулттай, кирилл, латин галиг, англи (заримд нь холимог)
хувилбартай, хариулах боломжгүй ба хоёрдмол асуулттай eval set бүтээж, gold бүрийг 9 өгөгдлийн сан дээр баталгаажуулав.
50 тохиргоог «нэг удаад нэг өөрчлөлт» зарчмаар (нийлмэл тохиргоог тусад нь тэмдэглэж) төлөвлөж, үнэлгээний дамжлагын
зөв байдлыг oracle run, аюулгүй байдлын 84 халдлага, 63 тест, кодын
хоёр удаагийн бие даасан хяналтаар баталгаажуулав. Мөн өмнөх ажлын системтэй ижил нөхцөлд харьцуулах орчин бүрдүүлэв.
\ifnum\resultsavailable=1
Үр дүнгийн дүгнэлтийг Ерөнхий дүгнэлтэд нэгтгэв.
\else
Загварын оноо, судалгааны асуултын (RQ1--RQ5) хариултыг серверийн туршилтын дараа энэ бүлэгт нэмнэ.
\fi
